\section{Εισαγωγή}\label{section:intro}

Στην παρούσα αναφορά εξετάζεται και σχεδιάζεται ένας απλός low-dropout (LDO) ρυθμιστής με pMOS στοιχείο διέλευσης (pass element). Το σύστημα θα πρέπει να λειτουργεί με τάση τροφοδοσίας $V_{\symup{in}}=1.2\unit{\volt}$ και στην έξοδό του να παρέχει σταθερή τάση $V_{\symup{out}}=0.8\unit{\volt}$ σε φορτίο σταθερού ρεύματος, $I_{\symup{load}}=1\unit{\milli\ampere}$.\par
Το LDO αποτελείται από έναν ενισχυτή σφάλματος (error amplifier - EA), ένα pMOS transistor ως στοιχείο διέλευσης και ένα δίκτυο ανάδρασης $R_{f}-R_{1}$.\par
Η σχεδίαση του συστήματος γίνεται στην τεχνολογία GPDK045. Λόγω της τάσης τροφοδοσίας, $1.2\unit{\unit\volt}$, επιλέχθηκαν τα transistor \texttt{nmos2v} και \texttt{pMOS2v} με ονομαστική τάση $1.8\unit{\volt}$ \cite{gpdk045ref}.

\begin{figure}[h!]
	\centering
	\begin{tikzpicture}
	% Paths, nodes and wires:
	\node[op amp] at (2, 4){};
	\node[vcc](N1) at (0, 4.539){} node[anchor=south] at (N1.text){$V_{\symup{ref}}$};
	\draw (0.81, 3.51) -| (0, 2);
	\draw (0, 2) to[american resistor, l={$R_{1}$}] (0, -0);
	\draw (1.5, 2) to[american resistor, l={$R_{f}$}] (3.5, 2);
	\node[vcc](N2) at (1.917, 4.539){} node[anchor=south] at (N2.text){$V_{\symup{in}}$};
	\node[vcc](N3) at (5, 4.539){} node[anchor=south] at (N3.text){$V_{\symup{in}}$};
	\draw (3.19, 4) -- (4.02, 4);
	\draw (0, 2) -- (1.5, 2);
	\draw (3.5, 2) -- (6, 2);
	\node[tlground] at (5, -0){};
	\node[tlground] at (0, -0){};
	\node[circ, xscale=0.75, yscale=0.75] at (0, 2){};
	\node[circ, xscale=0.75, yscale=0.75] at (5, 2){};
	\node[ocirc](N4) at (6, 2){} node[anchor=west] at (N4.text){$V_{\symup{out}}$};
	\draw (5, 2) -| (5, 2.961);
	\node[pigfete, solderdot](passElement) at (5, 3.73){} node[anchor=north east] at ([xshift=-0.27cm, yshift=-0.27cm]passElement.text){$\symup{M}_{\symup{pe}}$};
	\draw[-latex] (5, 3) -- (5, 2.5);
	\node[shape=rectangle, minimum width=0.465cm, minimum height=0.465cm](N5) at (5.25, 2.75){} node[anchor=center] at (N5.text){$I_{d}$};
	\draw[-latex] (5.25, 3.333) -- (5.25, 4.167);
	\node[shape=rectangle, minimum width=0.465cm, minimum height=0.465cm](N6) at (5.583, 3.667){} node[anchor=center] at (N6.text){$V_{sd}$};
	\draw[-latex] (4.25, 2) -- (3.75, 2);
	\node[shape=rectangle, minimum width=0.465cm, minimum height=0.465cm](N7) at (4, 1.70){} node[anchor=center] at (N7.text){$I_{f}$};
	\node[ground] at (1.917, 3.461){};
	\draw (0.81, 4.49) -| (0, 4.539);
	\draw (5, 4.5) -| (5, 4.539);
	\draw (5, 2) to[american current source, l={$I_{\symup{load}}$}] (5, -0);
\end{tikzpicture}
	\caption{Σχηματικό του ζητούμενου LDO με pMOS στοιχείο διέλευσης, $\symup{M_{pe}}$, και φορτίο σταθερού ρεύματος, $I_{\symup{load}}$. Είναι $V_{\symup{in}}=1.2\unit{\volt}$, $V_{\symup{out}}=0.8\unit{\volt}$, $V_{\symup{ref}}=0.6\unit{\volt}$ και $I_{\symup{load}}=1\unit{\milli\ampere}$.}
	\label{fig:LDO:schematic}
\end{figure}

\subsection{Αρχή Λειτουργίας}

Ο ενισχυτής σφάλματος συγκρίνει την τιμή της τάσης στην μη αναστρέφουσα είσοδο με την τάση αναφοράς, στην αναστρέφουσα είσοδο. Λόγω του διαιρέτη τάσης $R_f-R_1$, η τάση στην μη αναστρέφουσα είσοδο είναι
\begin{equation*}
	V_{+}=\frac{R_1}{R_f + R_1}V_{\symup{out}}.
\end{equation*}

Στην περίπτωση που είναι $V_{\symup{ref}}=V_{-}\gneqq V_{+}$ θα πρέπει η έξοδος του ενισχυτή να εφαρμόσει στην πύλη του στοιχείου διέλευσης $\symup{M_{pe}}$ κατάλληλο δυναμικό ώστε να άγει και η πτώση τάσης μεταξύ πηγής και εκροής να είναι $V_{sd}=0.4\unit{\volt}$. Εάν $V_{\symup{ref}}=V_{-}\lneqq V_{+}$ το στοιχείο διέλευσης δεν χρειάζεται να παρέχει ρεύμα στο φορτίο, κι επομένως, να άγει.\par
Οι τιμές των αντιστάσεων του δικτύου ανάδρασης επιλέχθηκαν έτσι, ώστε για $V_{\symup{out}}=0.8\unit{\volt}$ η τάση στην μη αναστρέφουσα είσοδο να είναι ίση με την τάση αναφοράς, $0.6\unit{\volt}$ και επιπλέον, όταν το στοιχείο διέλευσης άγει, το μεγαλύτερο μέρος του ρεύματος $I_{ds}$ να αποδίδεται στο φορτίο και το ρεύμα στο δίκτυο ανάδρασης, $I_f$, να είναι πολύ μικρό. Έτσι, για την $R_f$ επιλέχθηκε η τιμή των $47\unit{\kilo\ohm}$ και προέκυψε η $R_1=141\unit{\kilo\ohm}$.\par
Η $V_{sd}$ καλείται «dropout voltage» \cite{ChenPower}. Ο συντελεστής απόδοσης ισχύος, $\eta$, για το συγκεκριμένο σύστημα --όπως ορίζεται στο \cite{ChenPower}-- είναι $$\eta=\frac{V_{\symup{out}}}{V_{\symup{in}}}=\frac{0.8}{1.2}\unit{\volt\per{\volt}}=0.\overbar{6}$$

Ο εν λόγω LDO regulator ανήκει στους αναλογικούς LDO (A-LDO \cite{ChenPower}) και συγκεκριμένα στην υποκατηγορία των LDO χωρίς πυκνωτή παράλληλα στο φορτίο (\textsl{C-free} \cite{ChenPower}). Ο πυκνωτής παράλληλα στο φορτίο βελτιώνει τη συχνοτική απόκριση του συστήματος αναδιανέμοντας τους πόλους που προκύπτουν στην έξοδο του ενισχυτή και στο στοιχείο διέλευσης \cite{ChenPower} και βελτιώνει την απόκριση του συστήματος σε απότομες αλλαγές του φορτίου \cite{ChenPower}. Ωστόσο, η χωρητικότητά του κυμαίνεται στην τάξη των $\unit{\nano\farad}$ \cite{ChenPower} κι επομένως τοποθετείται εκτός του ολοκληρωμένου κυκλώματος \cite{ChenPower}.\par
Στις C-free αρχιτεκτονικές χρησιμοποιείται χωρητικότητα Miller μεταξύ των ενισχυτών και της εξόδου του συστήματος για ευστάθεια \cite{ChenPower}. Στην περίπτωση που το στοιχείο διέλευσης έχει μεγάλο λόγο $S=\nicefrac{L}{W}$, η παρασιτική χωρητικότητα πύλης-εκροής, $C_{gd}$, λειτουργεί σαν χωρητικότητα Miller \cite{ChenPower} εξαλείφοντας την ανάγκη για χρήση διακριτού πυκνωτή.\par

\subsection{Δομή της Αναφοράς}

Στην ενότητα \ref{section:EA} παρουσιάζεται η διαδικασία σχεδίασης του ενισχυτή σφάλματος (ενός σταδίου), σχετικές προσομοιώσεις, χαρακτηριστικά του ενισχυτή και το φυσικό του σχέδιο.\par
Στην ενότητα \ref{section:LDO} παρουσιάζεται η διαδικασία σχεδίασης του LDO και οι σχετικές προσομοιώσεις. Επιπλέον, προσομοιώνεται η συμπεριφορά του συστήματος με μεταβολή της τάσης τροφοδοσίας $\pm10\%$: $1.08\unit{\volt}\leqq V_{\symup{in}}\leqq 1.32\unit{\volt}$, και με μεταβολή του ρεύματος του φορτίου $\pm10\%$: $0.9\unit{\milli\ampere}\leqq I_{\symup{load}}\leqq 1.1\unit{\milli\ampere}$. Γίνονται προσομοιώσεις για TT, SS και FF corners. Τέλος, παρατίθεται και σχολιάζεται το φυσικό σχέδιο του LDO.\par
Στο παράρτημα \ref{appendix:layout} παρατίθενται το φυσικό σχέδιο του ενισχυτή σφάλματος και το φυσικό σχέδιο του LDO με μαύρο φόντο προκειμένου να είναι περισσότερο ευανάγνωστα.\par